\newpage
% Оформляем заголовок приложения согласно СТП БГУИР
\begin{flushright}
    {\bfseries ПРИЛОЖЕНИЕ Ж} \\
    (обязательное)
\end{flushright}

\begin{center}
    {\bfseries ОПИСАНИЕ ПРОГРАММЫ} \\
    \vspace{0.2cm}
    БГУИР ДП 1-40 01 01 001 13 01
\end{center}

\addcontentsline{toc}{section}{Приложение Ж. Описание программы}

\section*{Ж.1 Общие сведения}

Программное обеспечение «Система мониторинга параметров окружающей среды» (далее — Программа) предназначено для сбора данных с сенсоров, их первичной обработки и передачи на удаленный сервер по протоколу MQTT.

Программа написана на языке C++ с использованием фреймворка Arduino/ESP-IDF. Для разработки использовалась среда Visual Studio Code с расширением PlatformIO.

\section*{Ж.2 Функциональные возможности}

Программа обеспечивает выполнение следующих функций:
\begin{itemize}
    \item инициализация интерфейсов I2C и Wi-Fi;
    \item циклическое чтение данных из регистров датчика SHT31;
    \item программная фильтрация «шумов» измерений (метод скользящего среднего);
    \item контроль уровня заряда аккумуляторной батареи;
    \item управление энергопотреблением (переход в режим Deep Sleep).
\end{itemize}

\section*{Ж.3 Логическая структура}

Программа построена по модульному принципу. Основными структурными единицами являются:
\begin{enumerate}
    \item \textbf{Main module (main.cpp)} --- диспетчер задач, инициализация системы.
    \item \textbf{Network Manager} --- модуль управления Wi-Fi соединением и переподключением.
    \item \textbf{Sensor Driver} --- низкоуровневый драйвер для работы с шиной I2C.
    \item \textbf{MQTT Client} --- модуль формирования JSON-пакетов и взаимодействия с брокером.
\end{enumerate}



\section*{Ж.4 Технические средства}

Для функционирования Программы требуются следующие технические средства:
\begin{itemize}
    \item микроконтроллер семейства ESP32 (RAM 520 КБ, Flash 4 МБ);
    \item датчик температуры и влажности SHT3x;
    \item точка доступа Wi-Fi с выходом в интернет;
    \item MQTT-брокер (например, Mosquitto).
\end{itemize}

\section*{Ж.5 Вызов и загрузка}

Загрузка Программы в память микроконтроллера осуществляется через интерфейс UART с использованием утилиты \texttt{esptool.py}. После подачи питания или сброса Программа запускается автоматически из области памяти Flash.

\vfill
\begin{center}
    Минск 2026
\end{center}
\newpage